\SourceLeaf{335}

Sind Werthe von $\mathfrak{B}, \mathfrak{C}$ vorgegeben, die diesen Bedingungen, oder einer von ihnen, nicht genügen, so steht fest, dass diese Aufgabe eine physische Unmöglichkeit implicirt.

Umgekehrt, genügen die aufgegebenen Werthe den Bedingungen, so ist die Aufgabe möglich, d. i. die beiden Hülfskreise treffen sich in einem Punkte, der für beide Kreise im geeigneten Abschnitte liegt. Sind nemlich 1) die Zeichen jener Sinus $+\;+\;+$, so kann jener Durchschnitt nicht ausserhalb des Hauptkreises liegen, wo die ersten beiden Sinus negativ, der dritte positiv sein müsste; die beiden ersten Bedingungen würden nemlich erfordern, dass $\mathfrak{B}+180^0$, $\mathfrak{C}+180^0$ an die Stelle von $\mathfrak{B}, \mathfrak{C}$ treten, was aber wieder der dritten Bedingung widerspricht. Also liegt jener Durchschnitt innerhalb des Hauptkreises, wo eben die geeigneten Abschnitte liegen.

Auf ähnliche Weise wird 2) für die Zeichen $-\;-\;+$ bewiesen, dass der Durchschnitt nur ausserhalb liegen kann, wo dann die geeigneten Abschnitte liegen.

\begin{center}
\begin{tabular}{c@{\hspace{2em}}cc}
Sind die Zeichen & \multicolumn{2}{c}{so liegt der Durchschnitt} \\
 & im geeigneten Abschnitt & im ungeeigneten \\
$+\;+\;+$ & Kreis auf $a$ & [Kreis auf] $b, c$ \\
$+\;-\;+$ & $b$ & $a, c$ \\
$+\;-\;-$ & $c$ & $a, b$ \\
$-\;-\;-$ & $a$ & $b, c$ \\
$-\;+\;+$ & $c$ & $a, b$ \\
$-\;+\;-$ & $b$ & $a, c$ \\
\end{tabular}
\end{center}

Die elementarische Construction für das \textsc{Pothenot}sche Problem besteht darin, dass man Punkte in der durch (0) gehenden gegen 10 normalen Geraden bestimmt. Von einem solchen sind die Coordinaten polar relativ gegen (0):
\[ \text{Distanz } \frac{(12)}{\sin(02-01)}, \qquad \text{Richtung } 12+90^0-(02-01). \]

Eine andere neue Constructionsart ist die, dass man zunächst nicht den Punkt (0) sucht, sondern einen Punkt $(*)$, der in der Geraden 10 liegt und in der Distanz $(1*)=\frac{1}{(10)}$. Dieser Punkt wird durch Schnitte gerader Linien bestimmt, eine weniger, als von (0) aus Objecte geschnitten sind. Eine dieser
